\documentclass[12pt,a4paper]{exam}
\usepackage[T1]{fontenc}
\usepackage[utf8]{inputenc}
\usepackage[provide=*,lithuanian]{babel}
\usepackage{lmodern}
\usepackage[left=1.5cm,right=1.5cm,top=2.4cm,bottom=1.5cm]{geometry}
\usepackage{amsmath,amssymb,tasks,array}
\pagestyle{headandfoot}
\runningheadrule
\firstpageheader{Mokytojo variantas}{\shortstack{Kontrolinis darbas:\\Daugianariai}}{Šviesos variantas}
\runningheader{Daugianariai}{}{Šviesos variantas}
\firstpagefooter{}{\thepage}{}
\runningfooter{}{\thepage}{}
\pointsdroppedatright
\bracketedpoints
\pointname{ tašk.}
\newcommand{\atsakymas}{\par\medskip\textbf{Atsakymas:}\ \framebox[9cm]{\rule{0pt}{0.7cm}}}
\begin{document}

\begin{center}\Large\textbf{Sprendimai ir vertinimo kriterijai}\end{center}
\section*{Taškų paskirstymas ir sunkumas}
Sunkumas vertinamas mokiniui, kuris jau mokėsi neigiamųjų laipsnių, daugianarių tapatybių ir pilnojo kvadrato išskyrimo.
\begin{center}\begin{tabular}{|c|c|p{10cm}|}\hline
Užduotis & Taškai & Sunkumas ir pagrindimas\\\hline
1 & 6 & Lengva--vidutinė: šeši atskiri klasifikavimo atvejai, svarbu atskirti kintamąjį nuo iracionaliosios konstantos.\\\hline
2 & 5 & Vidutinė--sunki: keli neigiamieji laipsniai, laipsnių dalyba ir vienanario sąvoka.\\\hline
3(a) & 7 & Sunki: ilgas suprastinimas, pilnojo kvadrato išskyrimas ir maksimumo pagrindimas.\\\hline
3(b) & 3 & Vidutinė: didelio rodiklio lyginumas ir parametrų panaudojimas.\\\hline
4 & 5 & Vidutinė: koeficientų palyginimas ir tapatybės patikra.\\\hline
5 & 4 & Vidutinė: trijų pilnųjų kvadratų išskyrimas ir minimumo pasiekiamumas.\\\hline
Iš viso & 30 & \\\hline
\end{tabular}\end{center}
\section*{1 užduotis (6 taškai)}
\begin{center}
\begin{tabular}{|c|c|p{11cm}|}\hline
Dalis & Atsakymas & Paaiškinimas\\\hline
(a) & R, S & Vardiklis $\sqrt7$ yra nenulinė konstanta.\\\hline
(b) & R, T & $5x^{-3}=5/x^3$, todėl kintamasis yra vardiklyje; $x\ne0$.\\\hline
(c) & R, T & Kintamasis yra vardiklyje $x+2$; $\sqrt3$ yra konstanta; $x\ne-2$.\\\hline
(d) & I & Kintamasis yra po šaknimi; $x\leq7/2$.\\\hline
(e) & R, S & $16^{-5}$ ir $\sqrt{10}$ yra konstantos.\\\hline
(f) & R, T & Kintamasis yra vardiklyje $x^2$; $x\ne0$.\\\hline
\end{tabular}\end{center}
\textbf{Vertinimas:} po 1 tašką už kiekvieną tikslų raidžių rinkinį. Paaiškinimų ir apibrėžimo sričių mokinio sąlyga nereikalauja.
\section*{2 užduotis (5 taškai)}
Duotas reiškinys
$$\frac{(-2x^{11}y^{-3})^8\cdot(16x^{-2}y^5)^{-6}}{(32x^8y^{-9})^5}.$$
Pradinio reiškinio apibrėžimo sąlygos: $x\ne0$, $y\ne0$.
\begin{align*}
(-2x^{11}y^{-3})^8&=2^8x^{88}y^{-24},\\
(16x^{-2}y^5)^{-6}&=2^{-24}x^{12}y^{-30},\\
(32x^8y^{-9})^5&=2^{25}x^{40}y^{-45}.
\end{align*}
Todėl
\begin{align*}
\frac{2^8x^{88}y^{-24}\cdot2^{-24}x^{12}y^{-30}}{2^{25}x^{40}y^{-45}}
&=2^{8-24-25}x^{88+12-40}y^{-24-30+45}\\
&=2^{-41}x^{60}y^{-9}=\boxed{\frac{x^{60}}{2^{41}y^9}}.
\end{align*}
Pagal mokyklinę vienanario sampratą šis reiškinys \textbf{nėra vienanaris}, nes kintamojo $y$ laipsnio rodiklis yra neigiamas. Vienanario kintamųjų laipsnių rodikliai turi būti neneigiami sveikieji skaičiai.
\par\textbf{Vertinimas:} po 1 tašką už kiekvieną iš trijų teisingai pakeltų laipsniu daugiklių (3 tšk.); 1 taškas už galutinį supaprastinimą; 1 taškas už išvadą ir jos pagrindimą. Apibrėžimo sritis atskirai nevertinama.
\section*{3 užduoties (a) dalis (7 taškai)}
\begin{align*}
(2x-1)^2(x+1)&=(4x^2-4x+1)(x+1)=4x^3-3x+1,\\
(x-1)^3&=x^3-3x^2+3x-1,\\
(x+2)(x-1)(x+3)&=(x^2+x-2)(x+3)=x^3+4x^2+x-6.
\end{align*}
Taigi
\begin{align*}
E(x)&=4x^3-3x+1-(x^3-3x^2+3x-1)\\
&\quad-3(x^3+4x^2+x-6)+8x^2-19\\
&=-x^2-9x+1\\
&=-\left(x^2+9x+\frac{81}{4}\right)+\frac{81}{4}+1\\
&=\boxed{-\left(x+\frac92\right)^2+\frac{85}{4}}.
\end{align*}
Vadinasi, $a=-1$, $b=9/2$, $c=85/4$.
Kadangi $(x+9/2)^2\geq0$, tai $E(x)\leq85/4$.
Lygybė pasiekiama, kai $x=-9/2$.
\par\textbf{Atsakymas:} didžiausia reikšmė yra $85/4=21{,}25$, ji gaunama, kai $x=-9/2$. Mažiausios reikšmės nėra, nes $E(x)\to-\infty$, kai $|x|\to\infty$.
\par\textbf{Vertinimas:} po 1 tašką už tris išskleidimus (3 tšk.); 1 taškas už $-x^2-9x+1$; 1 taškas už pilnojo kvadrato išskyrimą; 1 taškas už maksimumo reikšmę ir pagrindimą; 1 taškas už $x=-9/2$.
\section*{3 užduoties (b) dalis (3 taškai)}
Pažymėkime $N=1\overbrace{0\ldots0}^{666\text{ nulių}}2$. Šis skaičius yra lyginis, nes baigiasi skaitmeniu 2. Todėl
$$a^N=(-1)^N=1,\qquad 21{,}25:c=\frac{85}{4}:\frac{85}{4}=1.$$
Gauname
$$b(a^N+21{,}25:c)=\frac92(1+1)=\boxed9.$$
\textbf{Vertinimas:} 1 taškas už laipsnio reikšmę, remiantis rodiklio lyginumu; 1 taškas už dalmens reikšmę; 1 taškas už galutinį atsakymą.
\section*{4 užduotis (5 taškai)}
Išskleidžiame kairiąją pusę:
$$(ax+a)(3x^2-bx+c)=3ax^3+a(3-b)x^2+a(c-b)x+ac.$$
Palyginę koeficientus, gauname
$$\begin{cases}3a=-6,\\a(3-b)=4,\\a(c-b)=2,\\ac=-8.\end{cases}$$
Iš pirmos lygybės $a=-2$. Tuomet
$$-2(3-b)=4\ \Longrightarrow\ b=5,\qquad -2c=-8\ \Longrightarrow\ c=4.$$
Patikriname likusią lygybę: $a(c-b)=-2(4-5)=2$. Visi koeficientai sutampa.
\par\textbf{Atsakymas:} $\boxed{a=-2,\ b=5,\ c=4,\ d\in\mathbb R}$.
\par\textbf{Vertinimas:} 1 taškas už bendrąjį išskleidimą; po 1 tašką už $a$, $b$, $c$ radimą (3 tšk.); 1 taškas už likusios koeficientų lygybės arba visos tapatybės patikrą.
\par\textit{Pastaba dėl sąlygos:} parametras $d$ pateiktoje lygybėje nedalyvauja, todėl jis gali būti bet kuris realusis skaičius. Jo paminėjimas laikomas sąlygos netikslumu, už jo aptarimą atskiras taškas neskiriamas.
\section*{5 užduotis (4 taškai)}
\begin{align*}
4x^2+12x&=(2x+3)^2-9,\\
9y^2-12y&=(3y-2)^2-4,\\
z^2-10z&=(z-5)^2-25.
\end{align*}
Todėl
\begin{align*}
&4x^2+9y^2+z^2+12x-12y-10z+40\\
&\qquad=(2x+3)^2+(3y-2)^2+(z-5)^2+2\geq2.
\end{align*}
Visi trys kvadratai vienu metu lygūs nuliui, kai
$$x=-\frac32,\qquad y=\frac23,\qquad z=5.$$
\textbf{Atsakymas:} mažiausia reikšmė yra $\boxed2$.
\par\textbf{Vertinimas:} po 1 tašką už kiekvieną teisingai išskirtą pilnąjį kvadratą (3 tšk.); 1 taškas už minimumą ir jo pasiekiamumo pagrindimą.
\section*{Bendros vertinimo nuostatos}
Lygiaverčiai teisingi sprendimo būdai vertinami taip pat. Už tą pačią klaidą pakartotinai taškai neatimami: tolesni metodiškai teisingi žingsniai vertinami pagal mokinio gautą rezultatą, jei klaida iš esmės nesupaprastino užduoties. Tai taikoma ir 3(b) daliai, naudojant 3(a) dalyje gautus parametrus.
\par\medskip

\begin{center}
\textbf{Vertinimo lentelė (iš viso 30 taškų)}\par\medskip
\setlength{\tabcolsep}{5pt}
\renewcommand{\arraystretch}{1.3}
\begin{tabular}{|c|*{10}{c|}}
\hline
Taškai & 0--3 & 4--6 & 7--9 & 10--12 & 13--15 & 16--18 & 19--22 & 23--25 & 26--28 & 29--30\\\hline
Pažymys & 1 & 2 & 3 & 4 & 5 & 6 & 7 & 8 & 9 & 10\\\hline
\end{tabular}
\end{center}
Procentinis rezultatas $p=100t/30$, kur $t$ -- surinkti taškai.
Kai $p<65$, skaičius $p/10$ apvalinamas į didesnę pusę;
kai $p\geq65$, jis apvalinamas iki artimiausio sveikojo skaičiaus
(pusė apvalinama į didesnę pusę). Mažiausias pažymys -- 1.
Procentinis rezultatas prieš skaičiuojant pažymį neapvalinamas.

\end{document}